\section{Contribución de vídeo}

En el contexto de producción en directo, se denomina \textit{contribución de vídeo}
al transporte de la señal audiovisual desde la fuente (cámara, ordenador, reproductor)
hasta el mezclador. Es el primer eslabón de la cadena y determina la calidad máxima
que el sistema puede procesar, así como la latencia de extremo a extremo.

\subsection{Interfaces físicas}

Las interfaces de contribución más extendidas en entornos profesionales son:

\begin{itemize}
  \item \textbf{SDI} (\textit{Serial Digital Interface}, SMPTE 259M/292M/424M):
  estándar de la industria broadcast para el transporte de vídeo sin comprimir sobre
  cable coaxial o fibra óptica. Soporta desde SD (270~Mbit/s) hasta 12G-SDI
  (11,88~Gbit/s para 4K). Proporciona latencia sub-imagen y es inherentemente síncrono.
  \item \textbf{HDMI}: interfaz digital de consumo, ampliamente disponible en cámaras
  y ordenadores. Aunque no está diseñada para largas distancias, su ubicuidad la
  convierte en la opción habitual en producciones de bajo coste cuando se emplean
  tarjetas de captura USB o PCIe.
  \item \textbf{Composite / Component}: interfaces analógicas en desuso, relevantes
  únicamente en instalaciones heredadas.
\end{itemize}

\subsection{Contribución sobre IP}

La tendencia actual en producción audiovisual es reemplazar el cableado dedicado por
transporte sobre redes Ethernet estándar mediante protocolos IP:

\begin{itemize}
  \item \textbf{NDI} (\textit{Network Device Interface}): protocolo propietario de
  NewTek que permite enviar vídeo sin comprimir o con compresión ligera sobre LAN con
  latencia de un fotograma. Es compatible con gran número de cámaras PTZ y software de
  producción.
  \item \textbf{SRT} (\textit{Secure Reliable Transport}): protocolo abierto desarrollado
  por Haivision que proporciona transmisión fiable de baja latencia sobre redes no
  confiables (Internet). Incorpora corrección de errores ARQ y cifrado AES-128/256,
  lo que lo hace idóneo para contribuciones remotas.
  \item \textbf{RTSP} (\textit{Real-Time Streaming Protocol}): protocolo de control
  empleado junto con RTP para la distribución de flujos multimedia en redes locales.
  Habitual en cámaras IP y sistemas de videovigilancia.
\end{itemize}

\subsection{Captura por software en Linux}

En sistemas Linux, la captura de señal de vídeo desde hardware se realiza
principalmente a través de \textbf{V4L2} (\textit{Video4Linux 2}), la API del
kernel que expone los dispositivos de captura (tarjetas capturadoras, webcams) como
ficheros de dispositivo (\texttt{/dev/video*}). GStreamer y FFmpeg interactúan
directamente con V4L2 para leer fotogramas y reencapsularlos en el formato que
voctocore espera (UYVY + PCM S16LE en Matroska/TCP).

Para entornos de desarrollo o pruebas sin hardware real, FFmpeg puede generar fuentes
sintéticas (barras de color, reloj en tiempo real) e inyectarlas directamente en
los puertos de entrada de voctocore, simulando el comportamiento de una cámara física.
